\section{Experiment}\label{sec:exp}

We evaluate CuFun on both synthetic and real-world datasets through
log-likelihood comparison, intensity estimation, and long-range dependency
analysis. Each experiment is designed to provide direct empirical evidence
for the theoretical guarantees established in Section~\ref{sec:thoe}.

% ------------------------------------------------------------------
\subsection{Datasets}
% ------------------------------------------------------------------

\paragraph{Synthetic Datasets.}

\noindent\textbf{Hawkes Processes.}
The Hawkes process employs a conditional intensity function
\[
  \lambda(t\mid\mathcal{H}(t))
  = \mu + \sum_{j=1}^M \alpha_j\beta_j\exp\!\bigl\{-\beta_j(t-t_j)\bigr\}.
\]
\textit{Hawkes1} uses $M=1$, $\mu=0.2$, $\alpha_1=0.8$, $\beta_1=0.8$.
\textit{Hawkes2} uses $M=2$, $\mu=0.2$, $\alpha_1=0.4$, $\alpha_2=0.4$,
$\beta_1=1.0$, $\beta_2=20.0$, introducing a fast-decaying component
that challenges numerical integration methods.

\noindent\textbf{Renewal Processes.}
Inter-event intervals $\{\tau_i=t_{i+1}-t_i\}$ are i.i.d.\ with density
$p(\tau)$~\cite{omi2019fully}.
\textit{S-Renewal} draws $p(\tau)$ from a log-normal distribution
(mean $1.0$, std $6.0$).
\textit{NS-Renewal} introduces non-stationarity~\cite{lindqvist2003trend}
by rescaling a gamma-distributed stationary sequence
($\mu=1.0$, $\sigma=0.5$) via the trend $r(t)=0.99\sin(2\pi t/20000)+1$,
verifying model adaptability under distributional shift.

\paragraph{Real-world Datasets.}
We use eight public datasets spanning diverse domains:

\noindent\textbf{Stack Overflow}~\cite{du2016recurrent,grant2013encouraging}:
Sequential badge-earning events from the most active users on the
programming Q\&A platform.

\noindent\textbf{MIMIC2}\footnote{https://github.com/musically-ut/tf\_rmtpp}:
Anonymised ICU visit records, where each event corresponds to a hospital admission.

\noindent\textbf{Twitter}\footnote{https://github.com/musically-ut/tf\_rmtpp}:
Retweet sequences time-stamped relative to the original tweet's creation.

\noindent\textbf{Reddit}\footnote{\label{fn:jodie}https://github.com/srijankr/jodie/}~\cite{kumar2019predicting,shchur2019intensity}:
Submission events on active subreddits from prominent users.

\noindent\textbf{Wikipedia}\textsuperscript{\ref{fn:jodie}}:
Edit events on the most active Wikipedia pages by highly active editors.

\noindent\textbf{Music}~\cite{celma2010music}:
User-specific music listening events.

\noindent\textbf{Bookorder}\textsuperscript{\ref{fn:jodie}}:
High-frequency trading records (buy/sell) from the New York Stock Exchange,
with millisecond resolution — a stress test for models under high-frequency events.

\noindent\textbf{Mooc}\textsuperscript{\ref{fn:jodie}}:
Student interactions (video views, problem-solving) on an online platform.

\noindent\textbf{Yelp}\footnote{https://www.yelp.com/dataset/challenge}:
Reviews of the top 300 most popular restaurants in Toronto, forming
per-restaurant temporal review sequences.

% ------------------------------------------------------------------
\subsection{Baselines}
% ------------------------------------------------------------------

We compare against four representative baselines:

\noindent\textbf{Constant Model}~\cite{sill1997monotonic,omi2019fully}:
An RNN that assumes a constant (exponential) inter-event distribution
whose mean is determined by the hidden state.
Numerical integration is required at inference, providing a lower-bound
reference for integration-induced error.

\noindent\textbf{Exponential (Exp) Model}:
An RNN with an intensity function that varies exponentially with $\tau$.
Numerical integration is unavoidable for log-likelihood computation.

\noindent\textbf{RMTPP}~\cite{du2016recurrent}:
A recurrent model with intensity as a non-linear function of event history.
Log-likelihood requires numerical integration, incurring the rounding error
analysed in Theorem~\ref{thm:numerical_superiority}.

\noindent\textbf{FullyNN}~\cite{omi2019fully}:
A feedforward network modeling the \emph{cumulative intensity}
$\Lambda(\tau)=\int_0^\tau\lambda^*(s)\,ds$, from which the CDF is recovered
as $F^*=1-\exp(-\Lambda)$.
This construction also guarantees a valid CDF and avoids quadrature at training
time.
Unlike CuFun, however, the additional nonlinear layer $\exp(-\Lambda)$ between
the network output and the log-likelihood increases the condition number of the
computation (Theorem~\ref{thm:numerical_superiority}), and the indirect CDF
parameterization does not carry the structural monotonicity guarantees of the
MNN architecture (Theorems~\ref{thm:mnn_monotone} and~\ref{thm:mnn_bounded}).

% ------------------------------------------------------------------
\subsection{Implementation Details}
% ------------------------------------------------------------------

Each experiment is repeated ten times with different train/validation/test
splits, with sequence length capped at 128.
We minimise NLL using Adam~\cite{kingma2014adam} (lr $=10^{-3}$, batch
size 64, up to 3000 epochs, early stopping after 100 stagnant epochs).
All hyperparameters (including $L_2$ regularisation) are held constant
across models. All experiments are implemented in PyTorch.

\paragraph{Evaluation metric.}
We report the Negative Log-Likelihood (NLL):
\begin{equation}
  \mathrm{NLL}(\theta)
  = -\sum_{i=1}^n \log P(y_i\mid\theta).
\end{equation}
Lower NLL indicates a better-fitting model.
NLL directly measures the quality of
$p^*(\tau\mid\boldsymbol{h}) = dF/d\tau$,
confirmed by Theorem~\ref{thm:cdf_intensity_equiv}.

% ------------------------------------------------------------------
\subsection{NLL Comparison on Synthetic Datasets}
% ------------------------------------------------------------------

Figure~\ref{nll_synthetic} reports NLL on the four synthetic datasets.
Three consistent findings emerge:

\paragraph{(1) CuFun achieves the lowest NLL on all synthetic datasets.}
CuFun attains NLL $0.471$, $-0.027$, $0.286$, $0.508$ on Hawkes1,
Hawkes2, S-Renewal, NS-Renewal — all strictly below FullyNN
($0.509$, $0.010$, $0.317$, $0.519$).
This validates Theorem~\ref{thm:universal_approx}: the MNN can
approximate any valid conditional CDF to arbitrary precision.

\paragraph{(2) Integration-based models degrade on high-frequency data.}
On Hawkes2 ($\beta_2=20.0$), the Exp model shows the largest gap.
This matches Theorem~\ref{thm:numerical_superiority}: large $N$ (needed
to resolve rapid oscillations) inflates $\gamma_N \gg \gamma_{2LW}$,
and accumulated rounding error dominates.
CuFun's autodiff with $\mathcal{O}(LW)$ operations stays numerically stable.

\paragraph{(3) CuFun outperforms FullyNN despite both avoiding quadrature.}
The advantage stems from \emph{structural} validity:
CuFun's monotone MNN guarantees $F\in[0,1]$ and $dF/d\tau>0$
(Theorems~\ref{thm:mnn_monotone}--\ref{thm:mnn_bounded}), letting the
optimiser exploit the full feasible parameter space without constraint
projection.

\begin{figure*}[!htp]
\centering
\begin{subfigure}{.235\textwidth}
  \centering\includegraphics[width=\linewidth]{Figures/Hawkes1.pdf}
\end{subfigure}%
\begin{subfigure}{.235\textwidth}
  \centering\includegraphics[width=\linewidth]{Figures/Hawkes2.pdf}
\end{subfigure}%
\begin{subfigure}{.235\textwidth}
  \centering\includegraphics[width=\linewidth]{Figures/S-Renewal.pdf}
\end{subfigure}%
\begin{subfigure}{.235\textwidth}
  \centering\includegraphics[width=\linewidth]{Figures/NS-Renewal.pdf}
\end{subfigure}
\caption{NLL comparison on synthetic datasets. \textbf{Lower is better.}
All improvements of CuFun are \textbf{statistically significant}
(two-sided $t$-test, $p<0.05$).}
\label{nll_synthetic}
\end{figure*}

% ------------------------------------------------------------------
\subsection{NLL Comparison on Real-world Datasets}
% ------------------------------------------------------------------

Figure~\ref{nll_real} reports relative NLL across eight real-world datasets
(positive = CuFun is better). Four observations stand out:

\paragraph{(1) CuFun achieves consistent superiority across domains.}
CuFun leads on six of eight datasets.
The gains on Bookorder (NLL $-8.060$ vs.\ FullyNN $-7.532$) and
Mooc ($6.155$ vs.\ $6.749$) are most pronounced — both high-frequency
datasets confirming Theorem~\ref{thm:numerical_superiority}.

\paragraph{(2) CDF modeling improves stability for extreme events.}
On Bookorder (millisecond-resolution timestamps spanning orders of
magnitude), intensity-based models suffer ill-conditioning when
$\lambda^*$ is very large or small.
CuFun's sigmoid output is always in $(0,1)$
(Theorem~\ref{thm:mnn_bounded}), preventing catastrophic cancellation
in $\log(1-F(\tau))$~\cite{magdon1998neural}.

\paragraph{(3) On social-media datasets, all models converge.}
Twitter (NLL $10.2043$) and Stack Overflow ($14.431$) show similar NLL
across models.

\paragraph{(4) Medical and music domains confirm broad applicability.}
On MIMIC2 ($0.466$) and Music ($-2.759$), CuFun maintains competitive
performance. The smooth, valid-CDF output supports
Theorem~\ref{thm:consistency_rate}: the approximation-error component
converges at the minimax-optimal rate $O(n^{-2\alpha/(2\alpha+d)})$,
benefiting moderate-sample domains.

\begin{figure*}[!tp]
\centering
\begin{subfigure}{.225\textwidth}
  \centering\includegraphics[width=\linewidth]{Figures/StackOverflow.pdf}
\end{subfigure}%
\begin{subfigure}{.225\textwidth}
  \centering\includegraphics[width=\linewidth]{Figures/MIMIC2.pdf}
\end{subfigure}%
\begin{subfigure}{.225\textwidth}
  \centering\includegraphics[width=\linewidth]{Figures/Twitter.pdf}
\end{subfigure}%
\begin{subfigure}{.225\textwidth}
  \centering\includegraphics[width=\linewidth]{Figures/Reddit.pdf}
\end{subfigure}
\newline
\begin{subfigure}{.225\textwidth}
  \centering\includegraphics[width=\linewidth]{Figures/Wikipedia.pdf}
\end{subfigure}%
\begin{subfigure}{.225\textwidth}
  \centering\includegraphics[width=\linewidth]{Figures/Music.pdf}
\end{subfigure}%
\begin{subfigure}{.225\textwidth}
  \centering\includegraphics[width=\linewidth]{Figures/Bookorder.pdf}
\end{subfigure}%
\begin{subfigure}{.225\textwidth}
  \centering\includegraphics[width=\linewidth]{Figures/Mooc.pdf}
\end{subfigure}
\caption{NLL comparison on real-world datasets (relative to CuFun).
\textbf{Lower is better.} All improvements are \textbf{statistically
significant} ($t$-test, $p<0.05$) except Wikipedia and Music.}
\label{nll_real}
\end{figure*}

% ------------------------------------------------------------------
\subsection{Intensity Function Estimation}
% ------------------------------------------------------------------

\subsubsection{High-Frequency Modeling on Hawkes2}

\begin{figure}[t]
\centering
\includegraphics[width=1.0\linewidth]{Figures/intensity_compa.pdf}
\caption{Estimated conditional intensity functions on the Hawkes2 dataset
($\beta_2 = 20.0$, fast-decaying component). CuFun tracks the true
generating intensity faithfully across both slow and fast time scales.
Exp and FullyNN show systematic deviations at high-frequency peaks,
consistent with the condition-number analysis of
Theorem~\ref{thm:numerical_superiority}.}
\label{simu_haw2}
\end{figure}

The Hawkes2 dataset provides a canonical stress test with slow
($\beta_1=1.0$) and fast ($\beta_2=20.0$) components.
We compare CuFun, FullyNN, and Exp against the true generating model.

The Exp model fails to track rapid oscillations
(Theorem~\ref{thm:numerical_superiority}): $N$-point quadrature inflates
$\gamma_N$ as $\beta_2$ increases, while CuFun's $\mathcal{O}(LW)$
autodiff keeps $\gamma_{2LW}$ small.

FullyNN, despite avoiding quadrature, shows frequent spikes in the
high-frequency regime.
Without the structural guarantees of the MNN architecture, density values
near zero create large $\log p^*$ penalties and gradient instability.
CuFun's bounded sigmoid output eliminates this failure mode.

\subsubsection{Long-Range Modeling on Yelp}

\begin{figure}[t]
\centering
\includegraphics[width=1.0\linewidth]{Figures/intenfun.pdf}
\caption{Density predictions on the Yelp dataset emphasising long-range
temporal modeling. CuFun captures the heavy tail more faithfully,
consistent with Theorem~\ref{thm:universal_approx}: the MNN can
$\varepsilon$-approximate any $\alpha$-H\"{o}lder CDF on $[0,T]$,
including heavy-tailed distributions.}
\label{density_yelp}
\end{figure}

The Yelp dataset has a heavy-tailed inter-arrival distribution,
testing long-range temporal modeling.
CuFun's density estimate (Figure~\ref{density_yelp}) matches the
empirical tail more faithfully than competitors.
By Theorem~\ref{thm:universal_approx}, CuFun's MNN can approximate any
$\alpha$-H\"{o}lder continuous CDF on $[0,T]$ to within $\varepsilon$,
including heavy-tailed ones.
The bias-variance decomposition (Theorem~\ref{thm:consistency_rate})
then guarantees minimax-optimal convergence over the full support, with
particular accuracy in the tail region governing rare high-impact events.

% ------------------------------------------------------------------
\subsection{Ablation: Fusion Operation}
% ------------------------------------------------------------------

\begin{figure}[t]
\centering
\includegraphics[width=0.8\linewidth]{Figures/two_value_final.pdf}
\caption{Value comparison during training: addition vs.\ element-wise
product operations. Under addition, a magnitude disparity causes
$\boldsymbol{h}$ to approach zero. Under element-wise product,
$\boldsymbol{h}$ maintains consistent scale, acting as a gate that
modulates future-event predictions from accumulated history.}
\label{value_comparison}
\end{figure}

As discussed in Section~\ref{sec:thoe}, it is essential to compare
the outputs from the history encoder ($\boldsymbol{h}$) and the
time-interval network ($\tau$).
We conduct this ablation on the Reddit dataset, where differences are
most pronounced.
Figure~\ref{value_comparison} plots the per-epoch output magnitudes of
the two branches.
In the addition variant, a significant magnitude disparity emerges:
$\boldsymbol{h}$ approaches zero, indicating that the history signal
is dominated by the $\tau$ branch.
In the element-wise product variant, both branches maintain consistent
scale; the history branch acts as a gate that multiplicatively
modulates the time-based prediction.
This gating mechanism is the key structural difference that enables
CuFun to leverage long-range dependencies effectively.

% ------------------------------------------------------------------
\subsection{Computational Efficiency Analysis}
% ------------------------------------------------------------------

\begin{figure}[t]
\centering
\includegraphics[width=0.6\linewidth]{Figures/inference_time_comparison (1).pdf}
\caption{Inference time comparison on the Hawkes1 dataset.
Time is measured in milliseconds (ms) per sequence on a
\textbf{logarithmic scale}. \textbf{Lower is better.}
All baselines, including FullyNN, exhibit significantly higher
inference times than CuFun.}
\label{fig:inference_time}
\end{figure}

We measure average wall-clock inference time (ms per sequence) on the
Reddit test set, with all models in the same environment averaged over
multiple runs.

As shown in Figure~\ref{fig:inference_time}, CuFun requires only
$31.0$ ms per sequence, versus $\sim\!140$ ms for all baselines —
nearly a $5\times$ speedup.
This gap applies equally to intensity-based models (RMTPP, Exp, Const)
and to FullyNN.

The underlying reason is algorithmic.
RMTPP and similar models are slow due to iterative Monte Carlo
simulation (e.g., Ogata's thinning algorithm).
FullyNN is slow at inference because predicting the next-event time
requires numerical root-finding on $F^*(\tau)=1-\exp(-\Lambda(\tau))$,
which in turn involves evaluating a neurally-parameterized integral.
CuFun circumvents both costs entirely: prediction reduces to a single
CDF inversion $\tau^* = F^{*-1}(u)$ for uniform $u\sim\mathrm{Unif}[0,1]$,
which is solvable in $O(1)$ network forward passes via bisection on a
bounded sigmoid output.
This analysis demonstrates that CuFun uniquely breaks the
efficiency-expressiveness trade-off, delivering state-of-the-art
accuracy with practical inference performance for real-world temporal
event forecasting.
